\textbf{Ice pellets}

An interesting weather phenomenon can occur when the temperature profile in the atmosphere shows an inversion. The solid blue line in figure 1 shows such a temperature profile. The inversion occurs at heights between $1\ km$ a$nd\ 2\ km$. Under these conditions snow falling through the atmosphere (partially) melts in the warmer layer and (partially) freezes again before reaching the ground in the form of ``ice pellets''.
% NOTE: "and" is printed with an upright "a" and italic "nd" (typesetting accident in the source)

\begin{center}\includegraphics[width=0.59\linewidth]{figures/EuPhO_2019_Q1_fig1.png}\end{center}

\begin{center}Figura 1: Atmospheric temperature $T$ vs. height $h$ above the ground.\end{center}

Assume that a small, spherical ice droplet almost completely melts while falling through the atmospheric layer between $h_A$ and $h_B$ where the temperature is above freezing point.

a) Determine the mass fraction of the droplet that freezes before reaching the ground.

b) Find, as precisely as possible, the temperature of the droplet at ground level if there were no inversion and the temperature profile followed the dashed line below a height of $2\ km$.

Neglect evaporation, condensation and size changes of the droplet. Assume that water and ice have very high thermal conductivity and that the density of the atmosphere is constant with height. Use $c_{water}=4.2\ kJ\ kg^{-1}\ K^{-1}$

for the the specific heat of water and $c_{ice}=2.1\ kJ\ kg^{-1}\ K^{-1}$ for that of ice. The specific latent heat for the melting of ice is $L=334\ kJ\ kg^{-1}$.
